\documentclass[11pt,fontset=fandol]{ctexart}
\usepackage[a4paper,margin=23mm]{geometry}
\usepackage{amsmath,amssymb,mathtools,bm,mathrsfs,esint}
\usepackage{graphicx,longtable,booktabs,array,calc}
\usepackage[normalem]{ulem}
\usepackage{xcolor,hyperref}
\hypersetup{hidelinks}
\setcounter{secnumdepth}{-1}
\newcounter{none}
\providecommand{\tightlist}{\setlength{\itemsep}{0pt}\setlength{\parskip}{0pt}}
\providecommand{\passthrough}[1]{#1}
\usepackage{amssymb,mathrsfs,esint,bm,newunicodechar}
\allowdisplaybreaks
\setlength{\emergencystretch}{3em}
\providecommand{\llangle}{\langle\!\langle}
\providecommand{\rrangle}{\rangle\!\rangle}
\newunicodechar{∞}{\ensuremath{\infty}}
\newunicodechar{∈}{\ensuremath{\in}}
\newunicodechar{∑}{\ensuremath{\sum}}
\newunicodechar{∫}{\ensuremath{\int}}
\newunicodechar{≤}{\ensuremath{\leq}}
\newunicodechar{≥}{\ensuremath{\geq}}
\newunicodechar{≠}{\ensuremath{\neq}}
\newunicodechar{→}{\ensuremath{\to}}
\newunicodechar{⇒}{\ensuremath{\Rightarrow}}
\newunicodechar{α}{\ensuremath{\alpha}}
\newunicodechar{β}{\ensuremath{\beta}}
\newunicodechar{γ}{\ensuremath{\gamma}}
\newunicodechar{λ}{\ensuremath{\lambda}}
\newunicodechar{φ}{\ensuremath{\varphi}}
\newunicodechar{ε}{\ensuremath{\varepsilon}}
\newunicodechar{σ}{\ensuremath{\sigma}}
\newunicodechar{θ}{\ensuremath{\theta}}
\newunicodechar{δ}{\ensuremath{\delta}}
\newunicodechar{η}{\ensuremath{\eta}}
\newunicodechar{π}{\ensuremath{\pi}}
\newunicodechar{ξ}{\ensuremath{\xi}}
\newunicodechar{Σ}{\ensuremath{\Sigma}}
\newunicodechar{Δ}{\ensuremath{\Delta}}
\newunicodechar{∂}{\ensuremath{\partial}}
\newunicodechar{∀}{\ensuremath{\forall}}
\newunicodechar{∃}{\ensuremath{\exists}}
\newunicodechar{√}{\ensuremath{\surd}}
\newunicodechar{−}{\ensuremath{-}}
\newunicodechar{μ}{\ensuremath{\mu}}
\newunicodechar{ρ}{\ensuremath{\rho}}
\newunicodechar{τ}{\ensuremath{\tau}}
\newunicodechar{ω}{\ensuremath{\omega}}
\newunicodechar{Ω}{\ensuremath{\Omega}}
\newunicodechar{χ}{\ensuremath{\chi}}
\newunicodechar{ψ}{\ensuremath{\psi}}
\newunicodechar{Γ}{\ensuremath{\Gamma}}
\newunicodechar{Φ}{\ensuremath{\Phi}}
\newunicodechar{Λ}{\ensuremath{\Lambda}}
\newunicodechar{ℝ}{\ensuremath{\mathbb{R}}}
\newunicodechar{ℤ}{\ensuremath{\mathbb{Z}}}
\newunicodechar{ℕ}{\ensuremath{\mathbb{N}}}
\newunicodechar{△}{\ensuremath{\triangle}}
\newunicodechar{⊗}{\ensuremath{\otimes}}
\begin{document}
\section{2026年高等代数真题}\label{ux5e74ux9ad8ux7b49ux4ee3ux6570ux771fux9898}

来源：\texttt{HDU/601数学分析/2026/2026年杭州电子科技大学数学分析+高等代数真题(双科合卷).pdf}。原
PDF 页序：3, 4。

\section{3
杭州电子科技大学2026年高等代数}\label{ux676dux5ddeux7535ux5b50ux79d1ux6280ux5927ux5b662026ux5e74ux9ad8ux7b49ux4ee3ux6570}

一、设 \(r = \sqrt [ n ] { p } ~ ( p\) 为素数），
\(f ( x ) = a _ { 0 } + a _ { 1 } x + \cdot \cdot \cdot + a _ { n - 1 } x ^ { n - 1 }\)
为整系数多项式，解答下列问题：

(1)证明 \(x ^ { n } - p\) 在有理数域上不可约；

(2)求有理数 \([ a , b , c ,\) 使得
\({ \frac { 1 } { 1 + { \sqrt [ { 3 } ] { 3 } } + { \sqrt [ { 3 } ] { 9 } } } } = a + b { \sqrt [ { 3 } ] { 3 } } + c { \sqrt [ { 3 } ] { 9 } } ;\)

(3)证明 \({ \frac { 1 } { f ( r ) } } = { \frac { D ( r ) } { D } }\)
，其中：

\[
D = \left( \begin{array} { c c c c } { a _ { 0 } } & { a _ { 1 } } & { \ldots } & { a _ { n - 1 } } \\ { p a _ { n - 1 } } & { a _ { 0 } } & { \ldots } & { a _ { n - 2 } } \\ { p a _ { n - 2 } } & { p a _ { n - 1 } } & { \ldots } & { a _ { n - 3 } } \\ { \vdots } & { \vdots } & { \ddots } & { \vdots } \\ { p a _ { 1 } } & { p a _ { 2 } } & { \ldots } & { a _ { 0 } } \end{array} \right) , \quad D ( r ) = \left( \begin{array} { c c c c } { 1 } & { a _ { 1 } } & { \ldots } & { a _ { n - 1 } } \\ { p } & { a _ { 0 } } & { \ldots } & { a _ { n - 2 } } \\ { p ^ { 2 } } & { p a _ { n - 1 } } & { \ldots } & { a _ { n - 3 } } \\ { \vdots } & { \vdots } & { \ddots } & { \vdots } \\ { p ^ { n - 1 } } & { p a _ { 2 } } & { \ldots } & { a _ { 0 } } \end{array} \right)
\]

二、设 \(A _ { m \times n } , B _ { n \times m }\) 满足 \(A B A = A\)
，β为m维列向量，解答下列问题：

(1)证明 \(A X = \beta\) 有解的充要条件是 \(A B \beta = \beta\)

(2)若 \(A X = \beta\) 有解，证明其所有解为
\(\displaystyle X = B \beta + ( E _ { n } - B A ) \alpha\)
（α为任意n维列向量）。

三、设\(\varphi\)是线性空间V上的线性变换，
\(\lambda _ { 1 } , \lambda _ { 2 } , \ldots , \lambda _ { s }\)
为的特征值，
\(\alpha _ { 1 } , \alpha _ { 2 } , \ldots , \alpha _ { s }\)
为对应特征向量，W是V的-不变子空间，k为W的维数。证明：若
\(\lambda _ { 1 } \alpha _ { 1 } + \lambda _ { 2 } \alpha _ { 2 } + \dots + \lambda _ { s } \alpha _ { s } \in W , \ k \geq s ,\)

四、设 \(f , g\)
是n维线性空间V上的线性变换，f有n个不同的特征值，解答下列问题：

(1)证明f的特征子空间均为g的特征子空间的充要条件是 \(f g = g f\)

(2)若 \(f g = g f\) ，证明g是ε，
\(f , f ^ { 2 } , \ldots , f ^ { n - 1 }\) 的线性组合（ε为恒等变换）。

五、设R2×2为2阶实矩阵空间，定义线性变换 \(f ( X ) = A X - X A\)
\(X \in \mathbb { R } ^ { 2 \times 2 }\)
\(A = { \left[ \begin{array} { l l } { 1 } & { 1 } \\ { 1 } & { 2 } \end{array} \right] }\)
再定义 \(( X , Y ) = \mathrm { t r } ( X ^ { T } A Y )\)
（tr为矩阵迹)，解答下列问题：

(1)证明(·,)是R2×2上的内积；

(2)求f的值域在该内积下的一组标准正交基。

六、设二次型
\(f ( x _ { 1 } , x _ { 2 } , x _ { 3 } , x _ { 4 } ) = X ^ { T } A X , X = ( x _ { 1 } , x _ { 2 } , x _ { 3 } , x _ { 4 } ) ^ { T }\)
，矩阵A满足：A有特征值入＝2,且
\(\mathrm { r a n k } ( 2 E _ { 4 } - A ) = 1\)
\(A = { \left[ \begin{array} { l l l l } { 2 } & { t } & { 2 } & { - 2 } \\ { a } & { 0 } & { b } & { c } \\ { d } & { e } & { 0 } & { f } \\ { g } & { h } & { k } & { 4 } \end{array} \right] }\)
，解答下列问题：

(1)判断A是否相似于对角矩阵，相似则证明，不相似则举反例；

(2)当t=2时，求f在正交变换下的标准形

七、设A,B均为正定矩阵，存在可逆矩阵P使得
\(P ^ { T } A P = E _ { n } , \ P ^ { T } B P = D = \mathrm { d i a g } ( \lambda _ { 1 } , \ldots , \lambda _ { n } )\)

解答下列问题：

(1)证明 \(\lambda _ { 1 } , \ldots , \lambda _ { n }\) 是
\(B A ^ { - 1 }\) 的全部特征值；

(2)证明A－B正定的充要条件是 \(B A ^ { - 1 }\) 的每个特征值满足
\(| \lambda _ { i } | < 1\)

(3)若 \(A ^ { 2 } - B ^ { 2 }\) 正定，证明 \(A - B\) 正定。

八、(1)若 \(A ^ { k } \alpha \neq 0 , A ^ { k + 1 } \alpha = 0\) ，证明
\(\alpha , A \alpha , \ldots , A ^ { k } \alpha\) 线性无关；

(2)证明
\(\operatorname{rank} ( A ^ { n } ) = \operatorname{rank} ( A ^ { n + 1 } ) = \operatorname{rank} ( A ^ { n + 2 } ) = . . .\)
。
\end{document}
